% ============================================================================
% CAPÍTULO 5 — Git\index{Git} e GitHub\index{GitHub}
% Status: texto completo (rodada editorial 2)
% ============================================================================
\chapter{Git e GitHub}
\label{cap:git}

\begin{caixaconceito}
\textbf{Git} é o sistema de controle de versão mais usado do mundo — o
``máquina do tempo'' do seu código. \textbf{GitHub} é a plataforma que hospeda
repositórios Git e adiciona colaboração, revisão e automação em cima dele.
\end{caixaconceito}

Você escreveu seu primeiro site (capítulos~\ref{cap:html} a~\ref{cap:js}) e ele
funciona. Agora vem a pergunta que todo profissional faz antes de continuar:
\emph{e se eu quebrar algo?} E se eu quiser testar uma ideia sem estragar o
que está pronto? E se eu precisar explicar, daqui a três meses, por que
mudei essa linha? A resposta a essas perguntas é o Git — a ferramenta mais
usada do desenvolvimento, e a primeira que você deve dominar antes de
qualquer framework.

Ao final deste capítulo, você será capaz de:
\begin{objetivos}
  \item Explicar o que o Git resolve e como funciona o modelo de \emph{snapshots};
  \item Executar o fluxo diário: \texttt{clone}, \texttt{status}, \texttt{add}, \texttt{commit}, \texttt{push}, \texttt{pull};
  \item Trabalhar com branches, merges e resolução de conflitos;
  \item Escrever mensagens de commit\index{commit} profissionais e commits atômicos;
  \item Colaborar via GitHub: Issues, Pull Requests e code review;
  \item Entregar o projeto do capítulo: o repositório profissional do seu portfólio.
\end{objetivos}

\section{O problema que o Git resolve}

Sem controle de versão, evoluir um projeto significa copiar pastas
(\path{site-final-v2-ok-agora-sim.zip}). O Git registra o \emph{estado
completo} do projeto a cada \texttt{commit} — quem mudou o quê, quando e por
quê — e permite voltar no tempo, ramificar experimentos e fundir trabalho em
paralelo sem perder nada.

\begin{caixaconceito}
Diferente de backups (que guardam cópias), o Git guarda \emph{snapshots} do
projeto inteiro a cada commit, com metadados completos (autor, data,
mensagem, hash). E como cada snapshot referencia o anterior, o histórico vira
um \emph{grafo navegável}: você pode comparar, reverter, ramificar e fundir
qualquer ponto da história.
\end{caixaconceito}

\subsection{O modelo de três áreas}

O Git funciona com três áreas que você precisa distinguir desde o início:

\begin{itemize}[leftmargin=*, itemsep=0.4em]
  \item \textbf{Working directory}: os arquivos que você edita no editor;
  \item \textbf{Staging area} (\texttt{git add}): o que você \emph{selecionou} para o próximo commit;
  \item \textbf{Repositório} (\texttt{git commit}): os snapshots gravados no histórico.
\end{itemize}

A separação entre \emph{editar}, \emph{selecionar} e \emph{gravar} é o que
permite commits atômicos: você pode trabalhar em três arquivos e gravar cada
mudança separadamente, com uma mensagem clara para cada uma.

\section{O fluxo diário}

\begin{lstlisting}[style=livro, language=Bash,
  caption={O ciclo de vida de um commit.}, label={list:git-fluxo}]
git clone https://github.com/usuario/projeto.git   # 1. baixa o projeto
cd projeto
git status          # 2. o que mudou?
git add index.html  # 3. seleciona o que entra no commit
git commit -m "feat: adiciona seção de contato"    # 4. grava o snapshot
git push            # 5. envia para o GitHub
git pull            # 6. traz mudanças dos outros
\end{lstlisting}

\begin{caixadica}
Use \texttt{git status} com frequência — ele é o ``GPS'' do Git: mostra onde
você está, o que mudou e o que está pronto para commitar. O comando
\texttt{git log --oneline} mostra o histórico de forma compacta.
\end{caixadica}

\section{Branches: experimentar sem quebrar}

Uma \emph{branch} é um ponteiro para um ponto do histórico — um ``universo
paralelo'' do seu projeto:

\begin{lstlisting}[style=livro, language=Bash,
  caption={Criando, trocando e fundindo branches.}, label={list:git-branches}]
git branch nova-funcionalidade   # cria a branch
git checkout nova-funcionalidade # ou: git switch nova-funcionalidade
# ... trabalha e commita ...
git checkout main                # volta para a principal
git merge nova-funcionalidade    # traz as mudanças para a main
git branch -d nova-funcionalidade # remove a branch já fundida
\end{lstlisting}

\begin{caixaconceito}
A regra de ouro: \textbf{a branch principal fica sempre funcionando}. Todo
experimento, correção ou feature nova nasce em uma branch\index{branch} própria e só volta
para a principal depois de testado. É isso que permite que várias pessoas —
ou vários ``eus'' — trabalhem sem pisar um no outro.
\end{caixaconceito}

\subsection{Conflitos: o Git pede sua decisão}

Quando duas branches mudam as mesmas linhas do mesmo arquivo, o Git não sabe
qual versão vale — e \emph{pede que você decida}. Um conflito aparece assim:

\begin{lstlisting}[style=livro, language=Bash,
  caption={Conflito de merge marcado no arquivo.}, label={list:git-conflito}]
<<<<<<< HEAD
const TITULO = "SkillHub";
=======
const TITULO = "Marketplace de serviços";
>>>>>>> feature/novo-nome
\end{lstlisting}

Você edita o arquivo escolhendo a versão correta (ou uma combinação), remove
os marcadores \texttt{<<<<<<<}, \texttt{=======} e \texttt{>>>>>>>}, e
finaliza com \texttt{git add} + \texttt{git commit}. Conflitos não são
erros — são o Git pedindo sua opinião. Quanto menores e mais frequentes os
commits, mais raros e simples os conflitos.

\section{Commits profissionais}

\begin{itemize}[leftmargin=*, itemsep=0.4em]
  \item \textbf{Atômicos}: um commit = uma mudança lógica (não ``mais um monte de coisa'');
  \item \textbf{Conventional Commits}: \texttt{feat:}, \texttt{fix:}, \texttt{refactor:}, \texttt{docs:}, \texttt{test:} — legível por humanos e ferramentas;
  \item \textbf{Mensagem boa}: o \emph{porquê} importa mais que o \emph{o quê} (``fix: valida e-mail antes de salvar'' \textbf{>} ``correção'').
\end{itemize}

\begin{caixadica}
Um bom teste para a mensagem do commit: ``se eu voltar a este commit daqui a
6 meses, eu entendo por que ele existe?''. O histórico bem escrito é a
documentação mais honesta do projeto — e a primeira coisa que um sênior lê
num code review.
\end{caixadica}

\section{GitHub: colaboração e revisão}

O GitHub adiciona a camada social e de processo sobre o Git:

\begin{itemize}[leftmargin=*, itemsep=0.4em]
  \item \textbf{Issue}: um problema, tarefa ou ideia registrada e rastreável;
  \item \textbf{Pull Request (PR)}: proposta de mudança de uma branch para outra, com discussão e revisão;
  \item \textbf{Code review}: comentários por linha, aprovações e sugestões;
  \item \textbf{README}: a porta de entrada do repositório (você vai escrever um de nível profissional no projeto deste capítulo).
\end{itemize}

\begin{caixaconceito}
O fluxo \textbf{feature branch $\rightarrow$ PR $\rightarrow$ review
$\rightarrow$ merge} é o coração do trabalho em equipe: nada entra na
principal sem passar por revisão. No capítulo~\ref{cap:cicd} você vai ver
como o CI automatiza as checagens nesse mesmo fluxo.
\end{caixaconceito}

% ============================================================================
\section{Projeto do capítulo: o repositório profissional do portfólio}
\label{sec:projeto5}

\begin{caixaprojeto}
\textbf{Objetivo:} transformar o site do capítulo~1 (portfólio) em um
repositório GitHub profissional: histórico limpo, README de nível
profissional e fluxo de PR.

\textbf{Critérios de aceite:}
\begin{itemize}[leftmargin=*, itemsep=0.2em]
  \item Repositório criado com \texttt{README.md}, \texttt{.gitignore} e licença;
  \item Histórico com commits atômicos usando Conventional Commits;
  \item Pelo menos uma branch de feature com merge\index{merge} via Pull Request\index{pull request};
  \item README explicando: o que é, como rodar, estrutura, captura de tela;
  \item \texttt{.gitignore} correto (sem \texttt{node\_modules}, \texttt{.env}, arquivos gerados);
  \item Deploy publicado (GitHub Pages\index{GitHub Pages} ou Vercel) com link no README.
\end{itemize}
\end{caixaprojeto}

\subsection{Passo a passo}

\begin{passos}
  \item \textbf{Local}: inicie o Git no projeto (\texttt{git init}), crie o
  \texttt{.gitignore} e faça o primeiro commit da base;
  \item \textbf{Histórico}: refaça o histórico em commits atômicos
  (\texttt{feat:}, \texttt{style:}, \texttt{fix:}) — se precisar, use
  \texttt{git reset} para reorganizar antes de publicar;
  \item \textbf{Remoto}: crie o repositório no GitHub e faça o primeiro
  \texttt{push};
  \item \textbf{README}: escreva um README profissional (descrição, captura,
  como rodar, estrutura, link do deploy);
  \item \textbf{PR}: crie uma branch \texttt{feature/melhorias}, faça uma
  melhoria real no site, abra o PR, revise e faça o merge;
  \item \textbf{Deploy}: publique no GitHub Pages (ou Vercel) e adicione o
  link ao README;
  \item \textbf{Documentação}: atualize a seção de referência do repositório
  e deixe o \texttt{git log} limpo como prova de trabalho profissional.
\end{passos}

\section{Exercícios de fixação}

\begin{exercicios}
  \item Explique a diferença entre working directory, staging area e repositório.
  \item O que faz \texttt{git status} e por que usá-lo com frequência?
  \item Escreva uma mensagem de commit para ``corrigi o cálculo do total no carrinho'' usando Conventional Commits.
  \item O que acontece quando duas branches alteram as mesmas linhas? Como você resolve?
  \item Liste o fluxo completo para enviar uma mudança do seu computador ao GitHub.
\end{exercicios}

\section{Desafios}

\begin{desafios}
  \item \textbf{Desfazer sem medo}: pratique \texttt{git revert} (commit novo que desfaz), \texttt{git reset --soft} (volta ao staging) e \texttt{git stash} — e explique quando usar cada um.
  \item \textbf{Histórico interativo}: use \texttt{git rebase -i} para squashar 3 commits de ``ajuste'' em 1 commit limpo antes do merge.
  \item \textbf{Colaboração real}: convide um colega (ou um segundo usuário seu) para abrir um PR no seu repositório e resolva um conflito de verdade.
\end{desafios}

\section{Exercícios de nível sênior}

\begin{seniores}
  \item \textbf{Política de branches}: desenhe a estratégia de branching do SkillHub (main protegida, PRs obrigatórios, tags de release) e documente-a.
  \item \textbf{Revisão de código}: faça um code review real de um PR público no GitHub — comente por linha, priorize os achados e explique a lógica de priorização.
  \item \textbf{Git para incidentes}: simule ``produção quebrou no commit X'' — use \texttt{git bisect} para encontrar o commit culpado em um histórico de 30 commits.
  \item \textbf{Automação}: crie um script de hook (\texttt{pre-commit}) que roda lint + testes antes de cada commit e documente o ganho.
\end{seniores}

\section{Armadilhas comuns}

\begin{itemize}[leftmargin=*, itemsep=0.4em]
  \item Commitar segredos (\texttt{.env}, chaves) — o histórico guarda para sempre;
  \item Commits gigantes sem mensagem útil (``update'', ``fix'', ``aaa'');
  \item Trabalhar direto na \texttt{main} (a branch principal deve ficar sempre funcionando);
  \item \texttt{git add .} sem revisar (o staging é seleção, não despejo);
  \item Deletar branches com trabalho não fundido (\texttt{-D} destrói commits);
  \item Medo de conflitos (conflitos são normais — resolva com calma e leitura).
\end{itemize}

\section{Resumo do capítulo}

\begin{itemize}[leftmargin=*, itemsep=0.3em]
  \item Git guarda snapshots com histórico completo; GitHub adiciona colaboração e revisão;
  \item Fluxo: \texttt{status} $\rightarrow$ \texttt{add} $\rightarrow$ \texttt{commit} $\rightarrow$ \texttt{push/pull};
  \item Branches isolam experimentos; a \texttt{main} fica sempre funcionando;
  \item Commits atômicos com Conventional Commits são a assinatura do profissional;
  \item PR + code review é o coração do trabalho em equipe;
  \item Projeto entregue: repositório profissional publicado e documentado.
\end{itemize}

\section{Referências oficiais para aprofundar}

\begin{itemize}[leftmargin=*, itemsep=0.3em]
  \item Documentação oficial do Git: \url{https://git-scm.com/doc}
  \item GitHub Docs: \url{https://docs.github.com}
  \item Conventional Commits: \url{https://www.conventionalcommits.org}
  \item Git — Pro Git (livro oficial, grátis): \url{https://git-scm.com/book/pt-br/}
\end{itemize}
